\section{Actual local fields, Frobenius signs and their complete intersection lattices}
Keep the original four labels $(\mathrm M,\mathrm N,\mathrm T,\mathrm E)
=(p,x,y,z)$, the positive ES equation, and the original coefficients
$S=p+x+y+z$, $A=-1/S$, $f(r)=\prod_i(r-r_i)$ from ISR1.
Here $p\equiv1\pmod {12}$ is prime, $R=\mathbb Z_p$ and
$K=\mathbb Q_p$. The positivity and reciprocity arguments in
PES5--7 and OE1--8 give the two exhaustive actual strata OE9.
Their ES parametrization retains the human provenance
\cite[Section 2, Type I and Type II parametrizations]{ElsholtzTao2015};
the signed factor coordinates retain \cite{ES488,Tao}.
No new coordinate, sign identification or quadratic twist is made here.
Write
\[
 d_i=f'(r_i),\qquad D_i=A d_i=p^{2m_i+\eta_i}u_i,
 \qquad \eta_i\in\{0,1\},\quad u_i\in R^\times,\qquad
 E=R[r,\sigma]/(f,\sigma^2-Af').                    \tag{SLF1}
\]
The ordered parity basis of $E$ and the complete original receiving
maps are exactly those of ISR4, SGR1 and SGR8.

\subsection{Square classes with every original unit retained}
For an odd prime, a unit $u\in R^\times$ is a square in $R$ exactly
when its residue is a square. Here is the full construction used
whenever a scalar square root occurs below. Choose
$a_1\in\{1,\ldots,p-1\}$ with $a_1^2\equiv u\pmod p$. Given
$a_n^2\equiv u\pmod {p^n}$, choose the unique $t_n\pmod p$ satisfying
$2a_nt_n\equiv(u-a_n^2)/p^n\pmod p$, and put
$a_{n+1}=a_n+p^nt_n$. The sequence is Cauchy, its limit has square
exactly $u$, and the other root is its negative. For a rational unit,
every displayed residue calculation uses its denominator inverse
modulo the stated power of $p$. Thus the construction keeps the full
unit, rather than replacing it by its residue representative.

If $u$ is a residue nonsquare, $R[\theta]/(\theta^2-u)$ has the
degree-two residue field and is the full ring of integers by the
trace argument in ISR5. Its nontrivial automorphism is
$\theta\mapsto-\theta$. On the residue field this is the
$p$-power map: Euler's criterion gives
$\bar\theta^p=\bar\theta(\bar u)^{(p-1)/2}=-\bar\theta$.
For two nonsquare units $u,v$, the ratio $u/v$ has square residue,
so the same recurrence constructs an exact isomorphism between
their two quadratic fields. The square root of $pu$, in contrast,
has valuation $1/2$, so its quadratic field is ramified and cannot
equal an unramified quadratic field. These observations prove all
field degrees and intersections used below.

In Type I put $z=pZ$, with $p\nmid xyZ$. All four full units are
\[
\begin{aligned}
 u_{\mathrm M}&=A(1-Z)(p-x)(p-y),&
 u_{\mathrm E}&=A(Z-1)(pZ-x)(pZ-y),\\
 u_{\mathrm N}&=A(x-p)(x-y)(x-pZ),&
 u_{\mathrm T}&=A(y-p)(y-x)(y-pZ).
\end{aligned}                                                   \tag{SLF2}
\]
Here $D_{\mathrm M}=pu_{\mathrm M}$,
$D_{\mathrm E}=pu_{\mathrm E}$ and
$D_{\mathrm N}=u_{\mathrm N}$, $D_{\mathrm T}=u_{\mathrm T}$.
The literal derivative product proves SLF2. Using $Z\equiv1/4$
and $A\equiv-1/(x+y)$ gives, with no zero denominator by OE8,
\[
 (\bar u_{\mathrm M},\bar u_{\mathrm N},
   \bar u_{\mathrm T},\bar u_{\mathrm E})
 =\left(-\frac{3xy}{4(x+y)},
 -\frac{x^2(x-y)}{x+y},\frac{y^2(x-y)}{x+y},
 \frac{3xy}{4(x+y)}\right).                            \tag{SLF3}
\]
Let $k=(a+b)/d$ be the exact Type I parameter in OE1. Its identity
$xy/(x+y)=Z/k$ proves
$D_{\mathrm M}/(pk)\equiv-3/(16k^2)$ and
$D_{\mathrm E}/(pk)\equiv3/(16k^2)$.
Both are squares: $-1$ is a square because $p\equiv1\pmod4$,
and $-3=(2\zeta+1)^2$ for a nontrivial cube root $\zeta$ in
$\mathbb F_p$, which exists because $p\equiv1\pmod3$.
Also $D_{\mathrm T}/D_{\mathrm N}\equiv-(y/x)^2$.
The recurrence therefore gives the full units
\[
 a_{\mathrm M}^2=\frac{D_{\mathrm M}}{pk},\quad
 a_{\mathrm E}^2=\frac{D_{\mathrm E}}{pk},\quad
 a_{\mathrm T}^2=\frac{D_{\mathrm T}}{D_{\mathrm N}},\qquad
 \omega^2=pk,\quad\theta^2=D_{\mathrm N}.              \tag{SLF4}
\]
For each selected residue root, the branch maps are exactly
$\sigma_{\mathrm M}\mapsto a_{\mathrm M}\omega$,
$\sigma_{\mathrm E}\mapsto a_{\mathrm E}\omega$,
$\sigma_{\mathrm N}\mapsto\theta$ and
$\sigma_{\mathrm T}\mapsto a_{\mathrm T}\theta$;
the inverse maps divide by the same units. Both choices of every
state sign remain present.

Put $\chi=\left(\frac{(x-y)/(x+y)}p\right)$. If $\chi=1$,
the two unit branch algebras split and the joint splitting field
is $K(\omega)$, with Galois sign image
\[
 G_{\rm I}=\{(+,+,+,+),(-,+,+,-)\}.
 \quad (\chi=1)                                       \tag{SLF5}
\]
If $\chi=-1$, then $K(\theta)$ is unramified quadratic and disjoint
from the ramified quadratic $K(\omega)$. The joint field has degree
four and sign image
\[
 G_{\rm I}=\{(+,+,+,+),(-,+,+,-),(+,-,-,+),(-,-,-,-)\}.
 \quad (\chi=-1)                                      \tag{SLF6}
\]
Its inertia subgroup is generated by $(-,+,+,-)$.
The Frobenius in its unramified quotient has the lift $(+,-,-,+)$
which fixes the specified ramified subfield $K(\omega)$; its other
lift is $(-,-,-,-)$. In SLF5 the unramified quotient is trivial.
Thus no unique Frobenius element on a ramified field has been
asserted. The two inertia signs at $\mathrm M$ and $\mathrm E$ are
the same because SLF4 proves an isomorphism using actual units,
not merely an equality of derivative valuations.

Both unit square classes of $k$ occur, independently of the sign
$\chi$, as witnessed by the exact positive solutions
\[
\begin{array}{c|c|c|c}
 (p;x,y,z)&k&\left(\frac{k}{p}\right)&\chi\\\hline
 (13;4,18,468)&11&-1&+1\\
 (13;4,20,130)&3&+1&-1\\
 (13;5,10,130)&3&+1&+1\\
 (37;16,22,6512)&19&-1&-1.
\end{array}                                                    \tag{SLF18}
\]
The ES identity and the Type I formula for $k$ are verified by
cross multiplication, and Euler's criterion gives the two symbol
columns. A ratio $pk_1/(pk_2)$ is a square in $K$ precisely when
$k_1/k_2$ has square residue, by the recurrence above. Consequently
the two possible ramified quadratic classes both occur, while
SLF4 continues to retain the original $k$ and every lifted scalar.

\subsection{All four Type II Frobenius patterns actually occur}
In Type II write $(y,z)=(pY,pZ)$; PES6 proves that $1,Y,Z$ are
pairwise distinct modulo $p$, and $S\equiv x\ne0$. The full units are
\[
\begin{aligned}
 u_{\mathrm M}&=A(p-x)(1-Y)(1-Z),&
 u_{\mathrm N}&=A(x-p)(x-pY)(x-pZ),\\
 u_{\mathrm T}&=A(pY-x)(Y-1)(Y-Z),&
 u_{\mathrm E}&=A(pZ-x)(Z-1)(Z-Y).
\end{aligned}                                                   \tag{SLF7}
\]
Here $D_i=p^2u_i$ for $i\in C=\{\mathrm M,\mathrm T,\mathrm E\}$
and $D_{\mathrm N}=u_{\mathrm N}$. Consequently
\[
 (\bar u_{\mathrm M},\bar u_{\mathrm N},
   \bar u_{\mathrm T},\bar u_{\mathrm E})
 =((1-Y)(1-Z),-x^2,(Y-1)(Y-Z),(Z-1)(Z-Y)).             \tag{SLF8}
\]
The second entry is a nonzero square. Multiplying all four full
derivatives, rather than only their residues, proves
\[
 \prod_i u_i
   =\left(\frac{A^2\det V}{p^3}\right)^2,\qquad
 \prod_i\left(\frac{u_i}p\right)=1,
 \quad V=(r_i^j)_{0\le i,j\le3}.                      \tag{SLF9}
\]
Indeed $\prod_i f'(r_i)=(\det V)^2$ in degree four and the sum
of the four derivative valuations is six. Thus either all four
branches split, or precisely two labels in $C$ are unramified
quadratic. The full list, with actual positive witnesses, is
\[
\begin{array}{c|c|c|c}
 (\epsilon_{\mathrm M},\epsilon_{\mathrm N},
  \epsilon_{\mathrm T},\epsilon_{\mathrm E})
 &(p;x,y,z)&(Y,Z)&\text{unique positive label in }C\\\hline
 (+,+,+,+)&(61;16,366,2928)&(6,48)&\text{all three}\\
 (+,+,-,-)&(13;4,26,52)&(2,4)&\mathrm M\\
 (-,+,+,-)&(37;10,148,740)&(4,20)&\mathrm T\\
 (-,+,-,+)&(61;18,122,549)&(2,9)&\mathrm E.
\end{array}                                                    \tag{SLF10}
\]
For example the last three reduced cluster triples
$(\bar u_{\mathrm M},\bar u_{\mathrm T},\bar u_{\mathrm E})$
are respectively $(3,11,6)$ modulo $13$, $(20,26,8)$ modulo $37$,
and $(8,54,56)$ modulo $61$. Euler's criterion gives exactly the
three rows. The first row gives $(52,34,22)$ modulo $61$, all
squares. Direct cross multiplication verifies the ES equation in
each row. Thus the list is exhaustive for every positive witness,
and every listed case occurs; the first and second rows alone
would not be an exhaustive list.

There are no ramified branches in this type. In a nontrivial row,
choose one of its two negative labels $a$, put $\theta^2=u_a$, and
write $\tau_i=p^{-m_i}\sigma_i$, with
$(m_{\mathrm M},m_{\mathrm N},m_{\mathrm T},m_{\mathrm E})=(1,0,1,1)$.
For each negative label the recurrence gives
$c_i^2=u_i/u_a$, and $\tau_i\mapsto c_i\theta$ is the exact
field isomorphism; choose $c_a=1$. For each positive label it gives
$b_i^2=u_i$ and the full split-algebra isomorphism
$\tau_i\mapsto(b_i,-b_i)\in K\times K$.
The state coordinates therefore retain
$\sigma_i=\pm p^{m_i}c_i\theta$ or
$\sigma_i=\pm p^{m_i}b_i$, respectively. The joint splitting field
is unramified quadratic. Its Frobenius is exactly the nontrivial
row of SLF10, since its residue map sends $\bar\theta$ to
$-\bar\theta$. In the first row the joint field is $K$ and the
Frobenius action on all eight states is the identity.

\subsection{The actual Galois-stable part of the original integral order}
For any sign vector retain the literal rational matrices
\[
 B_\epsilon=V^{-1}\operatorname{diag}(\epsilon_i)V,
 \qquad\Sigma_\epsilon=\operatorname{diag}(I_4,B_\epsilon).
                                                               \tag{SLF11}
\]
SLF4 and the Type II field maps identify these matrices with the
local Galois action on all eight states by an exact evaluation map.
For $s^2=D_i$, put
$e_{i,s}=(v(r_i)^{\mathsf T},s v(r_i)^{\mathsf T})$, where
$v(r)=(1,r,r^2,r^3)^{\mathsf T}$. If $g$ is the local Galois
element with sign vector $\epsilon$, then
$e_{i,s}\Sigma_\epsilon=e_{i,\epsilon_i s}=g(e_{i,s})$.
Thus the assertion concerns the action on the evaluated state,
with its fixed root label and both factor signs. Every element of the Type I group is
constant on the only nontrivial residue cluster $\{\mathrm M,\mathrm E\}$.
The exact Chinese-remainder idempotents in SGR2 therefore prove
$\Sigma_\epsilon E=E$. The largest suborder of $E$ stable under
the actual local Galois group is consequently $E$ itself in Type I.
This differs from the smaller suborder stable under all sixteen
formal independent signs in SGR4.

In a nontrivial Type II row, let $a,b$ be its two negative labels
in $C$, and let $j$ be the unique positive label in $C$.
The fourth label is the positive unit root $\mathrm N$.
Let $\Sigma_C$ have negative signs at all of $C$ and positive sign
at $\mathrm N$. The resultant between
$f_C=\prod_{i\in C}(r-r_i)$ and $r-x$ is a unit, because all
$r_i-x$ are units. Solving the Bezout identity over $R$ gives the
exact idempotent of $C$, hence $\Sigma_C\in\operatorname{GL}_8(R)$.
Moreover $\Sigma_\epsilon=\Sigma_C\Sigma_j$. Its integrality on a
vector is therefore exactly the integrality of $\Sigma_j$ on that
vector. The full SGR2 calculation gives
\[
\begin{aligned}
 J_G:=E\cap\Sigma_\epsilon E
   &=\{R_0+\sigma Q:Q(r_j)\in d_jR\},\\
 J_G&=D_jR^8,\qquad
 D_j=\operatorname{diag}(I_4,N_j),\qquad
 N_j=\begin{pmatrix}
 d_j&-r_j&0&0\\0&1&-r_j&0\\0&0&1&-r_j\\0&0&0&1
 \end{pmatrix},\\
 E/J_G&\simeq R/(d_j),\qquad
 R_0+\sigma Q\longmapsto Q(r_j)\pmod {d_j}.
\end{aligned}                                                   \tag{SLF12}
\]
The three possible full entries, not only their valuations, are
\[
\begin{aligned}
 d_{\mathrm M}&=p^2(p-x)(1-Y)(1-Z),\\
 d_{\mathrm T}&=p^2(pY-x)(Y-1)(Y-Z),\\
 d_{\mathrm E}&=p^2(pZ-x)(Z-1)(Z-Y).
\end{aligned}                                                   \tag{SLF13}
\]
Every displayed multiplier after $p^2$ is a unit by PES6.
Thus the defect is exactly one cyclic factor of length two,
and the relative odd-block Smith exponents of the actual Frobenius
matrix are $(-2,0,0,2)$. Multiplication by the integral invertible
$\Sigma_C$ preserves those exponents. Since the actual group has
order two, its two factors $E$ and $\Sigma_\epsilon E$ are permuted
by the group. Their intersection is a stable subring. Any stable
suborder contained in $E$ lies in that intersection, proving its
maximality. In the split Type II row, $J_G=E$ and the defect is zero.

The two eigenspaces also give a complete integral basis of this
intersection. Put $g_-=(r-r_a)(r-r_b)$ and
$g_+=(r-r_j)(r-x)$. Then
\[
\begin{aligned}
 J_G&=R[r]_{<4}\ \oplus\
  \sigma\bigl(g_-R[r]_{<2}\oplus g_+R[r]_{<2}\bigr),\\
 E^{G}&=R[r]_{<4}\oplus\sigma g_-R[r]_{<2},\qquad
 E^{-}=\sigma g_+R[r]_{<2},\\
 L_j&=\bigl([g_-],[rg_-],[g_+],[rg_+]\bigr),\\
 \det L_j&=(r_a-r_j)(r_a-x)(r_b-r_j)(r_b-x).
\end{aligned}                                                   \tag{SLF14}
\]
Brackets are the four literal coefficient entries in increasing
degree; no unit is omitted. To prove the decomposition, for
$Q\in J_G$ both $(Q+B_\epsilon Q)/2$ and
$(Q-B_\epsilon Q)/2$ are integral. The first vanishes at $r_a,r_b$
and the second at $r_j,x$. Repeated monic division proves divisibility
by the corresponding $g_-$ or $g_+$ over $R$; it does not require
the differences within $C$ to be units. Their sum is $Q$ and their
generic eigenspaces are disjoint. Conversely the displayed products
have signs $+1$ and $-1$ under $B_\epsilon$, so they belong to the
intersection. Expanding the four coefficient columns gives the
displayed resultant determinant, whose valuation is two.
In particular $L_jR^4=N_jR^4$ with their exact generators.

\subsection{Full return to the original weighted conductor}
Let $C_{\rm aff}$ be the full eight-by-four coefficient matrix in
SGR5, whose four coordinate functions are exactly FS13, and retain
the original $T_{\mathcal A}\Phi=\Psi$. Since $-1$ has a square
root in $K$, the displayed $i$ in those affine functions can be
given either chosen local embedding; both choices are fixed by the
Galois groups calculated here. More explicitly the common finite
algebra is
$\mathscr E_{\mathbb Q(i)}=\mathbb Q(i)[r,\sigma]/(f,\sigma^2-Af')$,
with its ordered eight-element parity basis. The full rational
coefficient matrix $\Sigma_\epsilon$ is an automorphism of this
algebra, and $C_{\rm aff}$ consists of its four specified coordinate
functions. Each embedding $\mathbb Q(i)\to\mathbb Q_p$ gives its
local realization; the fixed embedding $\mathbb Q(i)\to\mathbb C$
gives its complex realization. Its algebraic state evaluations
extend these embeddings to splitting fields separately. No map
from $\mathbb Q_p$ to $\mathbb C$ occurs. In the following formula choose either complex root
$s$ of $D_i$ and evaluate every point coordinate over $\mathbb C$;
the rational matrix $\Sigma_\epsilon$ carries the sign pattern
already determined by the local calculation. It gives the complete
signed affine return
\[
 q_{i,\epsilon_i s}
   =C_{\rm aff}^{\mathsf T}\Sigma_\epsilon^{\mathsf T}e_{i,s}^{\mathsf T},
 \qquad
 T_{\mathcal A}\Phi C_{\rm aff}^{\mathsf T}
       \Sigma_\epsilon^{\mathsf T}e_{i,s}^{\mathsf T}
    =\Psi q_{i,\epsilon_i s}.                         \tag{SLF19}
\]
This is an equality for the original four-dimensional point
coordinates, obtained by evaluating their exact coordinate
functions. It does not introduce a linearization of the nonlinear
sign involution on affine coordinates.

All coefficients in SLF11--14 are rational. Define the same modules
over $R_{(p)}=\mathbb Z_{(p)}$ and complete only for the local-field
calculation. Complex base change uses $R_{(p)}\to\mathbb Q\to\mathbb C$,
exactly as in SGR15; it introduces no map $\mathbb Q_p\to\mathbb C$.
For $\mathcal F_b=F_b\oplus F_b$ and the original
$\Gamma_b=\operatorname{diag}(H_b,H_b)$, $b=U,W$, put
\[
 \mathscr S_{b,\epsilon}=\mathcal F_b\Sigma_\epsilon\mathcal F_b^{-1}.
 \quad
 (T_{\mathcal A}\oplus T_{\mathcal A})\mathscr S_{U,\epsilon}
 =\mathscr S_{W,\epsilon}(T_{\mathcal A}\oplus T_{\mathcal A}).
                                                               \tag{SLF15}
\]
Every original moment, center and Gamma mass remains in $H_b$,
as specified by the original weighted conductor source \cite{WCF}.
For the nontrivial Type II Frobenius and each Type I generator with
two negative signs, use its actual two-positive/two-negative label
partition in $K_b=(V^{-1})^*H_bV^{-1}$.
If the corresponding cross-Gram singular values are
$\rho_{b,1}\ge\rho_{b,2}\ge0$, then the full original-metric
singular values, by the proved SGR11 computation, are
\[
 (s_1,s_2,1,1,1,1,s_2^{-1},s_1^{-1}),\qquad
 s_a=\sqrt{\frac{1+\rho_{b,a}}{1-\rho_{b,a}}},\qquad
 \|\wedge^k\mathscr S_{b,\epsilon}\|=
 \begin{cases}
 1&k=0,8,\\s_1&k=1,7,\\s_1s_2&2\le k\le6.
 \end{cases}                                                   \tag{SLF16}
\]
In the degree-four Type I case the two generators have identical
singular values: they differ by $\operatorname{diag}(I_4,-I_4)$,
which is an isometry for the unchanged $\Gamma_b$. Their product
has all eight singular values equal to one. No integrality argument
has been substituted for unitarity of either generator.

For each nontrivial Type II row the complete intersection return is
\[
\begin{aligned}
 (T_{\mathcal A}\oplus T_{\mathcal A})\mathcal F_UD_j
   &=\mathcal F_WD_j,\qquad \Gamma_{b,j}=D_j^*\Gamma_bD_j,\\
 \det\Gamma_{b,j}&=|d_j|^2\det\Gamma_b,\\
 \Gamma_{W,j}^{-1}\Gamma_{U,j}
   &=D_j^{-1}\Gamma_W^{-1}\Gamma_UD_j.
\end{aligned}                                                   \tag{SLF17}
\]
These are exact matrix products with the full SLF13 entry.
They prove that the arithmetic cyclic defect $R/(d_j)$ does not
alter any of the original inverse singular values or inverse
exterior norms when both metrics are pulled back along the same
complete intersection. The result identifies every actual local
Frobenius sign action and its exact original-metric image; it does
not identify that finite local action with a Frobenius realization
of the original global arithmetic conductor.
